\newpage
\phantomsection
\label{prf:boundedness-image-boundedness}
\LRAProofFor{prop:boundedness-image-boundedness}

\begin{remark*}[Return]
\hyperref[prop:boundedness-image-boundedness]{Return to Proposition}
\end{remark*}

\begin{proposition*}[Boundedness Is Image Boundedness]
For $f:A\to P$, being bounded above is equivalent to the existence of
$p\in P$ with $\forall a\in A,\ f(a)\leq p$.
\end{proposition*}

\begin{proof}[Professional Standard Proof]
\LRAProofBodyStart
The image $f(A)$ is bounded above by $p$ exactly when every element of
$f(A)$ is at most $p$. Since the elements of $f(A)$ are precisely the values
$f(a)$ with $a\in A$, this is equivalent to
$\forall a\in A,\ f(a)\leq p$. The lower-bound case is dual.
\end{proof}

\begin{proof}[Detailed Learning Proof]
\LRAProofBodyStart
To bound a function is to bound all its output values. Saying this through the
image and saying this pointwise over the domain are two readings of the same
condition.
\end{proof}
\begin{remark*}[Proof structure]
\end{remark*}



\begin{dependencies}
\begin{itemize}
  \item \hyperref[prop:boundedness-image-boundedness]{Proof target}
\end{itemize}
\end{dependencies}

\clearpage

